\documentclass[11pt]{article}
\usepackage[T1]{fontenc}
\usepackage{lmodern,amsmath,amssymb,amsthm,booktabs,array}
\usepackage[margin=1in]{geometry}
\usepackage[colorlinks=true,linkcolor=black,urlcolor=blue,citecolor=black]{hyperref}
\usepackage{microtype}
\newtheorem{theorem}{Theorem}
\newtheorem{lemma}[theorem]{Lemma}
\newtheorem{corollary}[theorem]{Corollary}
\newcommand{\rev}{\operatorname{rev}}
\newcommand{\pal}{\operatorname{Pal}}
\newcommand{\N}{\mathbb N}
\newcommand{\Z}{\mathbb Z}
\title{Carry Words and Reconstruction in\ Decimal Reverse-and-Add}
\author{Formal research note}
\date{13 October 2026}
\begin{document}
\maketitle
\begin{abstract}
We prove that every complete carry word produced by one decimal
reverse-and-add operation has a same-width positive representative reaching
a palindrome after at most four operations. The uniform bound four is
optimal. We also prove that equal complete carry histories containing five
non-overflow steps determine the initial mirrored digit sums, and hence the
next integer. Using the endpoint carry word sharpens this to a uniform
horizon of eleven consecutive carry words. We further prove that, at fixed
width, the numerical reverse-and-add result determines the entire pair-sum
word. The decimal results are formalized in Lean 4 using its
standard library. These are auxiliary results: no decimal Lychrel number is
constructed, and the status of 196 is not settled. Priority and publication
novelty have not been established.
\end{abstract}

\section{Definitions and statements}
Let $T(n)=n+\rev(n)$, where $\rev$ reverses the ordinary decimal digits,
discarding leading zeros in the reversed string. Let $T^k$ denote iteration.
A positive integer is palindromic when it equals its reversal.
The decimal Lychrel existence question asks whether some positive integer
has no palindromic iterate. This question, including the case of 196, remains
unresolved in this work.

Write an $L$-digit positive integer as
\[
 n=\sum_{i=0}^{L-1}d_i10^i,\qquad
 0\leq d_i\leq9,\quad d_{L-1}>0.
\]
Digits are indexed from the units position. Define the symmetric pair sums
and complete carry word by
\[
 q_i=d_i+d_{L-1-i},\qquad
 c_0=0,\qquad c_{i+1}=\left\lfloor\frac{q_i+c_i}{10}\right\rfloor.
\]
Thus $0\leq q_i\leq18$ and $c_i\in\{0,1\}$. The word
$(c_0,\ldots,c_L)$ includes the leading overflow $c_L$. A \emph{carry class}
consists of all positive integers of a fixed width having the same complete
carry word in their first operation.

\begin{theorem}[Sharp terminating representatives]\label{thm:four}
For every positive integer $a$, there is a positive integer $n$ with the
same decimal width and the same complete first carry word as $a$, and an
integer $k$ with $1\leq k\leq4$, such that $T^k(n)$ is a palindrome.
The uniform bound four cannot be replaced by three.
\end{theorem}

In contrast, consecutive carry words can recover information absent from
one carry word.
\begin{theorem}[Reconstruction from carry history]\label{thm:history}
Suppose two trajectories start at the same decimal width, and their complete
carry words agree at every step before a horizon $t$. If at least five of
these steps have no overflow, their initial pair-sum words agree. Consequently
their first reverse-and-add values agree.
\end{theorem}
The conclusion need not identify the initial integer: 196 and 295 both
have pair sums $(7,18,7)$ and first value 887.

\section{Carry-preserving normalization}
For each position set $a_i=c_{i+1}$ and $b_i=c_{L-i}$. These are the outgoing
carries at reflected columns. Define
\begin{equation}\label{eq:norm}
 \widehat q_i=\begin{cases}
 \ell,&a_i=b_i=0,\\
 10,&a_i=b_i=1,\\
 9,&a_i\ne b_i,
 \end{cases}
 \qquad
 \ell=\begin{cases}2,&c_L=0,\\0,&c_L=1.\end{cases}
\end{equation}

\begin{lemma}\label{lem:norm}
The word $\widehat q$ has exactly the same complete carry word as $q$.
It is the pair-sum word of a positive $L$-digit integer.
\end{lemma}
\begin{proof}
For a fixed sum $q_i=q_{L-1-i}$, two different outgoing carries are possible
only when the sum is nine: sums at most eight always produce zero and sums
at least ten always produce one. In the unequal-carry case the original
sum is therefore nine, and the replacement leaves it unchanged. In the
equal-carry cases, either choice $\ell\in\{0,2\}$ forces carry zero, and ten
forces carry one. Induction from $c_0=0$ proves equality at every position,
including the final overflow.

Reflection exchanges $a_i$ and $b_i$, so $\widehat q$ is symmetric. Its entries
are between zero and eighteen. At an odd center, $a_i=b_i$, so the central
sum is even. Its outer sum is positive: without overflow it is two or nine;
with overflow it is nine or ten. Each mirrored pair can therefore be split
into two decimal digits, with a positive leading digit, and an odd center
can be divided by two. These choices realize $\widehat q$ canonically.
\end{proof}

\subsection{No overflow}
When $c_L=0$, the next integer still has width $L$. If both reflected outgoing
carries vanish, its reflected digits sum to at most $2+1+2+1=6$.
If both are one, their sum is at most two. If they differ, the normalized
sum is nine, the incoming carries differ, and the output digits are zero
and nine. Every new pair sum is consequently at most nine. The following
addition has no carries and produces a palindrome. For example,
\[
196\text{ and }250\text{ have carry word }(0,0,1,0),
\qquad 250\longmapsto302\longmapsto505.
\]

\subsection{Overflow and protected runs}
For $c_L=1$, use $\ell=0$. The normalized alphabet is $\{0,9,10\}$.
Its outer symbol is either ten or nine. The exact pair-sum update, denoted
$G$, has the following properties:
\begin{align}
 \widehat q_0=10&\quad\Longrightarrow\quad
         \text{every entry of }G^2(\widehat q)\text{ is at most }9,
         \label{eq:ten}\\
 \widehat q_0=9&\quad\Longrightarrow\quad
         \text{every entry of }G^3(\widehat q)\text{ is at most }9.
         \label{eq:nine}
\end{align}
Here $G$ includes the changed reflection alignment when the width increases.
The overflow schedules are forced: ten becomes outer sum one, whereas nine
with overflow becomes ten and then one. Thus the schedules used in
\eqref{eq:ten} and \eqref{eq:nine} are respectively $(1,0)$ and $(1,1,0)$.
The one-digit case is $5\mapsto10\mapsto11$.

Here is a direct proof using runs of nines. Extend every word by zeros
outside its finite support. A symbol below nine resets the outgoing carry
to zero; ten resets it to one; nine transmits it. Call a word \emph{protected}
if every maximal nonempty run of nines has adjacent reset symbols zero and
ten, in either order. Let $P_d$ be the protected words whose other symbols
belong to $\{0,\ldots,d,10\}$, for $0\leq d\leq4$.

The normalized word belongs to $P_0$. On a run of nines, its two directional
carries differ by the normalization definition and are constant through the
run. They equal the types of the nearest reset symbols at opposite ends.
Those reset types must differ, so the two endpoints are zero and ten.
Exterior zeros supply a reset at an original endpoint.

Write $O$ and $N$ for the overflowing and non-overflowing pair-sum updates.
Let $R_i$ and $S_i$ be the output digits from carrying through the zero-extended
word from left and right, respectively. Symmetry gives the local formulas
\[
 N(q)_i=R_i+S_i,\qquad O(q)_i=R_i+S_{i-1}.
\]
Here zero extension is applied to $q$ before computing the streams:
$R_L=S_{-1}=h$ and $R_{-1}=S_L=0$, where $h$ is the overflow.
The shifted formula includes the additional leading digit automatically.
For two consecutive non-nine symbols $x,y$, their combined output in $O$ is
\[
 \rho(x)+\rho(y),\qquad
 \rho(x)=x\ (x<9),\quad \rho(10)=1.
\]
Indeed, each incoming carry is the reset type of the opposite symbol,
so their summed output is $x+y-9([x=10]+[y=10])$.

For a protected run, the output positions from the first old nine through
the old right endpoint are exactly
\begin{align*}
 0\,9^r\,10&\longmapsto10\,9^{r-1}\,0,\\
 10\,9^r\,0&\longmapsto0\,9^{r-1}\,10,\qquad r\geq1.
\end{align*}
Here $9^r$ denotes a run of $r$ symbols, not an integer power. The two streams
through the run are zero and nine. At its ends the sums are $9+1=10$ and
$0+0=0$, in the appropriate order. These identities imply
\[
 O(P_0)\subseteq P_2,\qquad O(P_2)\subseteq P_4.
\]
No new nine arises from ordinary pairs, and surviving runs retain opposite
zero/ten endpoints. Moreover, every four in $O(P_2)$ comes from adjacent
symbols $2,2$. Neither two can neighbor a nine in a protected word, so both
neighbors of the resulting four are below nine.

Under $N$, a protected run of nines stays nine: the opposing incoming carries
are zero and one. A ten produces at most two. An ordinary symbol $x$ produces
$2x$ plus its two incoming carries. Hence every symbol of $N(P_2)$ is at most
nine. For $N(O(P_2))$, values at most three produce at most eight; each four
has low neighbors, so both incoming carries vanish and its output is eight.
Thus $N(O(P_2))$ is also entirely low.

Starting in $P_0$, the forced schedules now give
\[
 P_0\xrightarrow{O}P_2\xrightarrow{N}\text{low},
 \qquad
 P_0\xrightarrow{O}P_2\xrightarrow{O}P_4
       \xrightarrow{N}\text{low}.
\]
The second arrow chain uses the neighbor property just proved, rather than
an assertion about every word in $P_4$. One further addition produces a
palindrome. This proves \eqref{eq:ten}--\eqref{eq:nine} and the upper bound
in Theorem~\ref{thm:four}. An independent finite-certificate proof underlies
the Lean formalization; its exact transition rules appear in Appendix~\ref{app:certificate}.

\section{Why four is optimal}
Consider the carry word $(0,0,0,1,1,1)$ at width five. Symmetry and its carry
equations force
\[
q=(9,9,e,9,9),\qquad e\in\{10,12,14,16,18\}.
\]
Indeed, $q_0<10$ and $q_0+1\geq10$ imply $q_0=9$; the reflected next pair
gives $q_1=9$ in the same way. The middle sum is twice a digit and at least
ten. Every representative therefore has one of only five possible first
values. The first three positive-time iterates are:
\begin{center}
\begin{tabular}{rrrr}
\toprule
$e$&$T(n)$&$T^2(n)$&$T^3(n)$\\
\midrule
10&100099&1090100&1101001\\
12&100299&1092300&1125201\\
14&100499&1094500&1149401\\
16&100699&1096700&1173601\\
18&100899&1098900&1197801\\
\bottomrule
\end{tabular}
\end{center}
None is a palindrome. On the other hand, the representative 99500 satisfies
\[
99500\longmapsto100099\longmapsto1090100\longmapsto1101001
\longmapsto2102012,
\]
whose last value is a palindrome. This proves optimality over entire carry
classes, rather than merely showing that one chosen normalization needs four
steps. The five-case exhaustion and the displayed orbit are checked in Lean.

\section{Reconstruction from several time steps}
For a trajectory define the polynomial
\[
 Q_t(x)=\sum_{i=0}^{L_t-1}q_{t,i}x^i.
\]
Consider two trajectories with identical complete carry histories and the
same initial width. Their overflow bits agree, so their widths agree at every
subsequent time. Let $\Delta_t$ be the difference of their pair-sum
polynomials. Subtracting the exact column equations cancels every carry term:
\begin{equation}\label{eq:diff}
 \Delta_{t+1}(x)=
 \begin{cases}
 2\Delta_t(x),&\text{no overflow},\\
 (1+x)\Delta_t(x),&\text{overflow}.
 \end{cases}
\end{equation}
The second line includes both new outer positions and the shifted interior.
Thus, after $N$ non-overflow and $H$ overflow steps,
\[
 \Delta_t(x)=2^N(1+x)^H\Delta_0(x).
\]
If $\Delta_0\ne0$, let $m$ be its least nonzero exponent and $a$ the
corresponding integer coefficient. The coefficient at $x^m$ on the right
is $2^Na$, since $(1+x)^H$ has constant coefficient one. Its absolute value
is at least $2^N$. But every coefficient on the left lies between $-18$ and
$18$. When $N\geq5$, this contradicts $32>18$. Hence $\Delta_0=0$, proving
Theorem~\ref{thm:history}. The Lean proof follows the least differing
coefficient directly, avoiding a polynomial-library dependency.

\subsection{A uniform eleven-word horizon}
The five-step condition can be guaranteed at a fixed temporal horizon.
Let $a=q_0\in\{1,\ldots,18\}$ be the current outer pair sum, $a'$ the next,
and $h$ the overflow. The possible outer transitions are contained in
\[
\begin{array}{c|c|c}
a&h&a'\\\hline
1\leq a\leq8&0&2a\text{ or }2a+1\\
9&0&18\\
9&1&10\\
10\leq a\leq18&1&a-9.
\end{array}
\]
Define a potential by listing its values at $a=1,\ldots,18$:
\[
(\psi(1),\ldots,\psi(18))=
(0,0,0,1,0,1,1,1,2,1,1,1,2,1,2,2,2,3).
\]
Each listed transition satisfies $2h-1\leq\psi(a)-\psi(a')$.
Summing through $t$ steps and using $0\leq\psi\leq3$ gives
$2H-t\leq3$. At $t=12$, this already implies $H\leq7$ and $N\geq5$.
An endpoint refinement gives the stronger bound below.

\begin{corollary}[Eleven consecutive carry words]\label{cor:eleven}
The initial decimal width and complete carry words at times
$0,1,\ldots,10$ determine the initial pair-sum word.
\end{corollary}
\begin{proof}
Equal incoming and outgoing carries at a column place both pair sums in
the same interval
\[
10c_{i+1}-c_i\ \leq\ q_i\ \leq\ 10c_{i+1}+9-c_i.
\]
Their difference therefore has absolute value at most nine. After ten
operations, the growth bound gives $H\leq6$, hence $N\geq4$. If an initial
least differing coefficient existed, its difference at time ten would have
absolute value at least $2^4=16$. The equal carry word at time ten bounds
this same difference by nine, a contradiction.
\end{proof}
This corollary, including the outer potential and its telescoping argument,
is also kernel-checked. No optimality is claimed for eleven.

\section{Uniqueness of the numerical predecessor word}
The numerical evaluation of a pair-sum word and the next pair-sum word must
be distinguished. Write
\[
 S(q)=\sum_{i=0}^{L-1}q_i10^i,
 \qquad G(q)=\text{pair-sum word of the canonical digits of }S(q).
\]
Thus $S(q)=T(n)$ when $q$ is the pair-sum word of $n$.

\begin{theorem}[Fixed-width uniqueness]\label{thm:inverse}
Let $q,r$ be symmetric words of the same length $L$, with entries between
zero and eighteen. If $S(q)=S(r)$, then $q=r$.
\end{theorem}
\begin{proof}
The two additions have identical output digits and identical final carries.
Let $e_i$ be the first carry minus the second carry at position $i$.
Then $e_0=e_L=0$, each $e_i\in\{-1,0,1\}$, and the column equations give
\[
 q_i-r_i=10e_{i+1}-e_i.
\]
Suppose inductively that $e_i=e_{L-i}=0$. Reflection symmetry of $q-r$
then gives
\[
 10e_{i+1}=-e_{L-1-i}.
\]
The right-hand side has absolute value at most one, so both sides vanish.
Induction from the two ends makes every carry difference zero. The column
equations now give $q_i=r_i$ at every position.
\end{proof}

In particular, two positive integers of the same width have the same first
reverse-and-add value if and only if their pair-sum words agree. Positivity
and the even-center condition are not needed in the algebraic theorem.
The same argument works in every base $b\geq2$ with entries at most $2b-2$;
the accompanying Lean theorem is decimal.

\begin{corollary}
A positive decimal target has at most two canonical pair-sum predecessor
words, one at its own width and one at one less than its width.
\end{corollary}
\begin{proof}
A reverse-and-add operation preserves width or increases it by one.
Theorem~\ref{thm:inverse} gives at most one word for each possible width.
\end{proof}

This bound cannot be iterated to count predecessors under $G$. To see why,
choose $x_1,\ldots,x_m$ independently from $\{1,10\}$. All $2^m$ words
\[
 (10,x_1,\ldots,x_m,x_m,\ldots,x_1,10)
\]
have the same $G$-image: outer entries one and every interior entry two.
Indeed, no entry equals nine, and an overflow transition combines adjacent
values of $\rho(x)=x-9\mathbf1_{x\geq10}$. Both $\rho(1)$ and $\rho(10)$
equal one. The numerical outputs $S(q)$ are distinct, as the theorem
requires; taking their pair sums loses information. An existence argument
based on a uniform bound of two predecessors for $G$ would therefore be
invalid.

\section{Formalization, scope, and reproducibility}
The self-contained source \texttt{Lychrel196Research.lean} uses
\texttt{import Std} and Lean 4.19.0. Its arithmetic definitions use ordinary
natural-number decimal reversal. Canonical digit representations include a
nonzero leading digit; trailing zeros are permitted. The numeric statements
therefore concern usual integers, not padded strings or an abstract dynamical
system lacking an arithmetic interpretation.

Principal declarations are:
\begin{itemize}
\item \texttt{CarryWordTermination.}\\
\texttt{every\_positive\_integer\_has\_fast\_same\_carry\_representative};
\item \texttt{CarryNormalizationSharpness.four\_is\_optimal};
\item \texttt{CarryHistory.iterPairs\_carry\_history\_injective};
\item \texttt{CarryHistory.eleven\_carry\_words\_determine\_pair\_sums};
\item \texttt{ReverseAddInverse.symmetric\_value\_injective};
\item \texttt{ReverseAddInverse.same\_width\_step\_eq\_iff\_pair\_sums}.
\end{itemize}
The first proves the unconditional numeric four-step theorem, the second its
sharpness, the third the five-non-overflow reconstruction theorem for
exact pair-sum orbits, and the fourth Corollary~\ref{cor:eleven}.
The last two prove Theorem~\ref{thm:inverse} and its equivalence for actual
positive decimal integers. The two-width corollary and exponential-fiber
example are proved in prose above, not formalized in Lean.
The accompanying build log records the complete
source hash, compiler result, and axiom audit. There are no \texttt{sorry}
proof holes, custom axioms, or \texttt{native\_decide} assertions. Standard
Lean axioms such as propositional extensionality, choice, and quotient
soundness may occur. Compile with
\begin{verbatim}
lean Lychrel196Research.lean
\end{verbatim}
The archive also contains the modular sources and a Python implementation
that regenerates the finite normalization graphs. Python and SAT-solver
results from unrelated construction searches are not premises of the
formal theorems above.

The normalization theorem rules out a universally sound Lychrel criterion
depending only on one current complete carry word: every such word also
occurs in a terminating trajectory. It does not rule out a criterion using
digit values, several time layers, or additional restrictions proved for
the particular orbit of 196. Reconstruction demonstrates precisely why
multi-time information escapes this obstruction.

None of these theorems constructs a nonterminating trajectory. The formal source
retains two unproved definitions, named \texttt{Decimal196Conjecture}
and \texttt{DecimalLychrelExistence}. The associated searches for
closed, palindrome-avoiding families have so far produced no witness.

\paragraph{Related work and priority.}
Carry descriptions of one-step palindrome formation appear in the
correspondence of Vaughn Suite hosted by p196.org~\cite{suite}. A recent
manuscript by Ando discusses carry asymmetry and a general-base one-step
characterization~\cite{ando}. Those subjects are distinct from the
carry-class normalization and history-reconstruction statements proved
here. The limited literature check does not establish that these statements
are new. This document is a reviewable research note, not a claim of a solved
decimal Lychrel conjecture or of established publication priority.

\begin{thebibliography}{9}
\bibitem{suite} Vaughn Suite, correspondence hosted by Wade Van Landingham,
\emph{196 and Other Lychrel Numbers},
\url{https://www.p196.org/html/papers5.html}.
Accessed 13 October 2026. The displayed no-carry inequalities in the
correspondence contain a strict-bound typo; the formal work here uses
the exact condition $q_i\leq9$.
\bibitem{ando} Ando, \emph{Carry Asymmetry, Palindrome Characterization,
and Conditional Non-Convergence in the Reverse-and-Add Process},
repository manuscript, February 2026,
\url{https://github.com/nkc-daiki/196-carry-asymmetry}.
Accessed 13 October 2026. Cited as related work, not as a premise of the proofs.
\end{thebibliography}
\appendix
\section{Exact finite certificate}\label{app:certificate}
For completeness, we specify the finite certificate proving these two
arbitrary-width statements. Its entire transition definition, explicit state
sets, closure checks, and arithmetic interpretation are in the accompanying
Lean source. It is a finite-state proof, not a search through integers below
a cutoff.

\paragraph{Normalized input.}
Read mirrored pairs from outside inward. State $(a,b)\in\{0,1\}^2$
records the low entering carry and high outgoing carry. Start at $(0,1)$.
The transitions are exactly
\[
\begin{array}{c|c|c}
\text{condition}&\text{emitted pair sum}&\text{next state}\\\hline
b=0&0&(0,b'),\quad b'\in\{0,1\}\\
b=1&10&(1,b'),\quad b'\in\{0,1\}\\
a\ne b&9&(a,b).
\end{array}
\]
At an even center accept if $a=b$. At an odd center the central sum is $10b$.
Restrict the first pair to ten or nine for the respective claim.

\paragraph{One exact time layer.}
Let
\[
 C(q,a)=\lfloor(q+a)/10\rfloor,\qquad D(q,a)=(q+a)\bmod10.
\]
For a prescribed overflow $h$, the state is $(a,b,u)$; $u$ is a delayed high
output digit when $h=1$. Begin with $(0,h,1)$ when $h=1$, and $(0,0,0)$
when $h=0$ (the last coordinate is then unused). On input $q$, choose
$b'\in\{0,1\}$ satisfying $C(q,b')=b$. Put
\[
 a'=C(q,a),\quad v=D(q,a),\quad w=D(q,b').
\]
Emit $v+w$ and move to $(a',b',0)$ when $h=0$; emit $v+u$ and move to
$(a',b',w)$ when $h=1$. The delayed digit accounts exactly for the shifted
mirror caused by overflow.

At an even input center require $a=b$; output no center when $h=0$, and
center $2u$ when $h=1$. At an odd input center $z$, require $C(z,a)=b$.
For $h=0$ output center $2D(z,a)$. For $h=1$ emit the final pair $D(z,a)+u$
and no center. A final pair is passed through all remaining time layers
before their center checks are performed.

\paragraph{Closure certificate.}
Compose the input automaton with the required time layers, retaining a
Boolean flag that becomes true if any final-layer pair exceeds nine. A
completion is unsafe if that flag is true or its final center exceeds nine.
The explicit closed state sets have the following sizes:
\begin{center}
\begin{tabular}{ccrrr}
\toprule
First pair&Schedule&States&Edges&Center completions\\
\midrule
10&$(1,0)$&37&89&11\\
9&$(1,1,0)$&159&391&25\\
\bottomrule
\end{tabular}
\end{center}
Every possible initial state belongs to its respective set, and every
successor of every state remains in that set. Every accepted
completion is safe. These finite propositions are proved by kernel reduction
in Lean; induction on the number of input pairs then covers every finite
width. The semantic bridge identifies the accepted transitions and center
conditions with the actual numeric operation. This proves
\eqref{eq:ten}--\eqref{eq:nine}. One further carry-free operation yields a
palindrome, proving the upper bound in Theorem~\ref{thm:four}.



\clearpage
\section{Extension to sufficiently large bases}
\label{app:bases}
This appendix is a mathematical extension of the direct proof. Its
base-parametric statement is not formalized in the accompanying Lean file;
the fully formalized specialization is decimal.

\begin{theorem}\label{thm:bases}
Let $b$ be even with $b\geq10$, or odd with $b\geq17$. Every complete
first carry word in base $b$ has a positive same-width representative
with a base-$b$ palindromic iterate at some strictly positive time
$1\leq k\leq4$. The bound is sharp among such positive times.
\end{theorem}
\begin{proof}
Put $s=1$ for even $b$ and $s=2$ for odd $b$, and set
$\alpha=b+s-1$. Thus $\alpha$ is even and $b-1\geq8s$.
Normalize equal outgoing carries to $2$ or $\alpha$ in the no-overflow
case, and to $0$ or $\alpha$ in the overflow case; leave unequal carries
at $b-1$. The same induction preserves all carries. The evenness of
$\alpha$ ensures a valid central digit. Without overflow, the next pair
sums are at most six, at most $2s$, or exactly $b-1$, so the following
operation is palindromic.

For overflow, enlarge the protected-run condition: each run of $b-1$ has
endpoint pair
\[
 (a,b+s-1-a),\qquad 0\leq a\leq s-1,
\]
in either orientation. Normalization starts with $(a,b+s-1-a)=(0,\alpha)$.
The overflow rule on such a run is
\[
 a\,(b-1)^r\,(b+s-1-a)
 \longmapsto (b+a)\,(b-1)^{r-1}\,(s-1-a).
\]
This preserves the permitted endpoint class. For odd bases the types
$(0,b+1)$ and $(1,b)$ exchange; retaining only the first type would be
incorrect.

For ordinary adjacent symbols use $\rho(x)=x$ below $b-1$, and
$\rho(x)=x-b+1\in\{1,\ldots,s\}$ for high symbols. The first overflowing
image has ordinary low symbols at most $2s$, and the second at most $4s$.
A maximal low symbol $4s$ comes only from adjacent $2s,2s$. Neither can
neighbor a neutral run, since allowed low endpoints are at most $s-1$.
Its two new neighbors are low, so the subsequent no-overflow output is
$8s$. Every smaller low symbol produces at most $8s$ as well; high symbols
produce at most $2s$, and protected neutral runs remain $b-1$.
All are below $b$.

The outer value $\alpha$ becomes $s$, forcing schedule $O,N$; the outer
value $b-1$ becomes $b$, then one, forcing $O,O,N$. The same argument
therefore gives a palindrome after at most four operations.

For sharpness, the five-digit carry word $(0,0,0,1,1,1)$ forces
$q=(b-1,b-1,b+z,b-1,b-1)$, where $0\leq z\leq b-2$ and $b+z$ is even.
Writing digit tuples least-significant first, its first three iterates are
\[
\begin{aligned}
 &(b-1,b-1,z,0,0,1),\\
 &(0,0,z+1,z,b-1,0,1),\\
 &(1,0,z,2z+1-b\kappa,z+\kappa,1,1),
 \quad \kappa=\lfloor(2z+1)/b\rfloor.
\end{aligned}
\]
The first two fail at their outer digits, and the third at its second
outer pair. At the least permitted $z=s-1$, the fourth iterate is
$2102012_b$ for even $b$, or $2126212_b$ for odd $b$, both palindromes.
\end{proof}
For odd bases in this theorem, every normalized pair sum is even, so splitting
each in half already gives a palindromic representative at time zero.
The sharpness assertion concerns strictly positive iteration times; it does
not assert that this representative is initially nonpalindromic.

\end{document}
